\section{卷积运算的数学细节}

\subsection{输出尺寸计算}

给定输入空间尺寸 $W$、滤波器大小 $K$、填充 $P$ 和步幅 $S$，输出尺寸为：

$$W_{out} = \left\lfloor \frac{W - K + 2P}{S} \right\rfloor + 1$$

\begin{figure}[H]
\centering
\includegraphics[width=0.95\textwidth]{slides-images/slide-044.jpg}
\caption{不同超参数对输出尺寸的影响示例\protect\footnotemark}
\end{figure}
\footnotetext{Slides 第45页。}

\subsection{填充（Padding）}

\begin{figure}[H]
\centering
\includegraphics[width=0.95\textwidth]{slides-images/slide-046.jpg}
\caption{零填充：在输入边缘添加零值，使得卷积后空间尺寸保持不变\protect\footnotemark}
\end{figure}
\footnotetext{Slides 第47页。}

不使用填充时，每次卷积都会使空间尺寸缩小 $K-1$ 个像素。这在深层网络中会导致特征图迅速缩小到 0。

\begin{importantbox}{保持空间尺寸的常用设置}
当 $K$ 为奇数时，设置 $P = \frac{K-1}{2}$、$S = 1$ 可以保持卷积前后的空间尺寸不变。这是最常见的卷积配置。例如：
\begin{itemize}
    \item $K = 3, P = 1, S = 1$
    \item $K = 5, P = 2, S = 1$
    \item $K = 7, P = 3, S = 1$
\end{itemize}
\end{importantbox}

\subsection{步幅（Stride）与下采样}

\begin{figure}[H]
\centering
\includegraphics[width=0.95\textwidth]{slides-images/slide-048.jpg}
\caption{步幅为 2 的卷积：滤波器每次跳过一个位置，有效地对特征图进行下采样\protect\footnotemark}
\end{figure}
\footnotetext{Slides 第49页。}

步幅大于 1 的卷积会\textbf{缩小}空间尺寸——这是在网络内部实现\textbf{下采样}的一种方式。下采样非常重要，因为它允许网络的感受野以指数速度增长。

\subsection{感受野（Receptive Field）}

\begin{figure}[H]
\centering
\includegraphics[width=0.95\textwidth]{slides-images/slide-050.jpg}
\caption{感受野的增长：单层 $3 \times 3$ 卷积的感受野为 3，三层 $3 \times 3$ 卷积的感受野增长到 7\protect\footnotemark}
\end{figure}
\footnotetext{Slides 第51页。}

\textbf{感受野}（receptive field）指的是：网络某一层的一个激活值，能``看到''原始输入图像中多大的区域。

\begin{importantbox}{感受野的增长规律}
\begin{itemize}
    \item \textbf{无下采样}（$S = 1$）：感受野\textbf{线性}增长。$L$ 层 $K \times K$ 卷积的感受野为 $L(K-1) + 1$。这意味着需要很多层才能覆盖大图像的全部区域。
    \item \textbf{有下采样}（$S = 2$）：感受野\textbf{指数}增长。较少的层即可覆盖整张图像。
\end{itemize}
最终分类决策需要聚合整张图像的信息，因此下采样对于构建实用的 CNN 至关重要。
\end{importantbox}

\subsection{1x1 卷积}

\begin{figure}[H]
\centering
\includegraphics[width=0.85\textwidth]{slides-images/slide-053.jpg}
\caption{$1 \times 1$ 卷积：在每个空间位置上对通道维度做全连接变换\protect\footnotemark}
\end{figure}
\footnotetext{Slides 第54页。}

$1 \times 1$ 卷积是一种特殊的卷积：空间上不做任何邻域操作，只在\textbf{通道维度}上做线性组合。它可以看作在每个空间位置独立应用的全连接层。主要用途：
\begin{itemize}
    \item 改变通道数（降维或升维）
    \item 增加非线性（配合激活函数）
    \item 减少计算量（先降维再做大卷积）
\end{itemize}

\subsection{卷积层的计算量分析}

\begin{figure}[H]
\centering
\includegraphics[width=0.95\textwidth]{slides-images/slide-055.jpg}
\caption{卷积层的计算量分析：FLOPs、参数量、内存占用\protect\footnotemark}
\end{figure}
\footnotetext{Slides 第56页。}

对于一个卷积层（$C_{in}$ 输入通道，$C_{out}$ 输出通道，$K \times K$ 滤波器，输出空间大小 $H' \times W'$）：

\begin{table}[H]
\centering
\begin{tabular}{ll}
\toprule
\textbf{指标} & \textbf{数值} \\
\midrule
参数量 & $C_{out} \times C_{in} \times K \times K$ \\
FLOPs & $C_{out} \times C_{in} \times K \times K \times H' \times W'$ \\
输出内存 & $C_{out} \times H' \times W'$ \\
\bottomrule
\end{tabular}
\caption{卷积层资源消耗}
\end{table}

\begin{knowledgebox}{CNN 中计算量和参数量的分布}
在典型的 CNN 中：
\begin{itemize}
    \item \textbf{前面的卷积层}：空间尺寸大、通道数少，计算量大但参数少
    \item \textbf{后面的卷积层}：空间尺寸小、通道数多，计算量相对小但参数多
    \item \textbf{全连接层}：参数量最多但计算量相对少
\end{itemize}
这种分布特性对网络设计和硬件优化都非常重要。
\end{knowledgebox}

\subsection{本章小结}

卷积运算的数学细节虽然看起来复杂，但核心思想简单：在局部区域做内积，然后滑动。填充保持空间尺寸，步幅实现下采样，感受野随深度增长。理解这些细节对于设计和调试 CNN 架构至关重要。

%% ========== Part 6: 池化层 ==========

